% ============================================================
%  Section 2, Segment 2A — Basic Music Theory / 基础乐理
%  原书页码 21-27
% ============================================================
\chapter{Understanding Basic Music Theory / 基础乐理入门}
\label{ch:theory}

\eng{In this section, I will provide an overview of music theory relevant to math rock and related music styles. We will explore fundamental principles of music theory, offering you a solid starting point to expand your knowledge. Before we dive in, you might have two essential questions:}

\chn{在本节中，我将概述与数学摇滚及相关音乐风格相关的乐理知识。我们将探索音乐理论的基本原理，为你提供拓展知识的坚实起点。在深入之前，你可能会有两个关键问题：}

\begin{itemize}[leftmargin=*]
  \item \eng{How Much Music Theory Knowledge is Needed to Learn, Write, and Play Math Rock?}
  \chn{学习、创作和演奏数学摇滚需要多少乐理知识？}
  \item \eng{What Music Notation Concepts Should I Know to Best Learn, Write, and Play Math Rock?}
  \chn{为最好地学习、创作和演奏数学摇滚，我应该掌握哪些乐谱概念？}
\end{itemize}

\eng{How Much Music Theory Knowledge is Needed to Learn, Write, and Play Math Rock?}

\chn{学习、创作和演奏数学摇滚需要多少乐理知识？}

\begin{itemize}[leftmargin=*]
  \item \eng{Key signatures (what sharps and flats are in each key: accidentals)}\chn{调号（每个调中有哪些升降号：临时记号）}
  \item \eng{Semitones and tones}\chn{半音与全音}
  \item \eng{Chords that are built from each note in a key (diatonic)}\chn{从调内每个音建立的和弦（调式内和弦）}
  \item \eng{Intervals (space between two notes)}\chn{音程（两个音之间的距离）}
  \item \eng{Chord progressions}\chn{和弦进行}
  \item \eng{Melody \& harmony}\chn{旋律与和声}
  \item \eng{Rhythm (patterns of strong and weak pulses)}\chn{节奏（强弱脉冲的模式）}
  \item \eng{Time signatures}\chn{拍号}
  \item \eng{Major \& minor scales}\chn{大调与小调音阶}
  \item \eng{Chord types - Triads, extended chords, sus chords, etc}\chn{和弦类型——三和弦、扩展和弦、挂留和弦等}
  \item \eng{Repetition and variation}\chn{重复与变化}
  \item \eng{Strumming patterns}\chn{扫弦模式}
\end{itemize}

\eng{What Music Notation Concepts Should I Know to Best Learn, Write, and Play Math Rock? (These are all related to reading music on a stave):}

\chn{为最好地学习、创作和演奏数学摇滚，我应该掌握哪些乐谱概念？（这些都与五线谱阅读有关）：}

\begin{itemize}[leftmargin=*]
  \item \eng{Time signatures (simple, compound, odd)}\chn{拍号（简单拍、复拍、奇数拍）}
  \item \eng{Key signatures}\chn{调号}
  \item \eng{Note divisions (whole, half, quarter, eighth, sixteenth, triplets, etc)}\chn{音符时值（全音、二分音、四分音、八分音、十六分音、三连音等）}
  \item \eng{Reading rhythm}\chn{阅读节奏}
  \item \eng{Rests}\chn{休止符}
  \item \eng{Ties and dots}\chn{连线与附点}
  \item \eng{Repetition signs}\chn{重复记号}
  \item \eng{Notation symbols for guitar: hammer-ons, pull-offs, slides, bends, staccato, tapping (left and right hand), - guitar pro is a great way to learn these notation symbols!}\chn{吉他记号：击音、勾音、滑音、弯音、断音、点弦（左右手）——Guitar Pro 是学习这些记号的好工具！}
  \item \eng{Dynamical markers}\chn{力度记号}
\end{itemize}

\eng{Learning these concepts is not as challenging as you might imagine, and the effort you invest in understanding them is highly rewarding. I recommend using this guide as a starting point and also obtaining a grade 1 music theory book, then gradually progressing through to grade 5. I myself have achieved a grade 5 in music theory, which is often suggested as the goal for those learning popular music styles. Beyond grade 5, music theory exams tend to focus more on classical music.}

\chn{学习这些概念并不像你想象的那么困难，投入理解它们的努力也会带来丰厚的回报。我建议把本指南作为起点，同时准备一本 1 级音乐理论书，然后逐步学习到 5 级。我自己已获得音乐理论 5 级，这通常被认为是学习流行音乐风格的目标。超过 5 级后，音乐理论考试通常会更多地关注古典音乐。}

\paragraph{Major and Minor}
\eng{Let's start with major and minor key signatures. Most modern math rock is quite melodic and usually uses chords and notes from one key signature (be this a conscious decision or not) and typically is comprised of harmony and melody you typically hear in popular music styles. This is somewhat different to earlier variations of math rock music, which is often referred to as being more dissonant in nature. I have linked an article below that was written by Nickk Hunter of Fecking Bahamas which clarifies these differences better than I can here.}

\chn{让我们从大调和小调调号开始。大多数现代数学摇滚相当富有旋律性，通常使用来自一个调号的和弦与音符（无论这是有意为之还是无意识中形成），并且通常包含类似流行音乐的和声与旋律。这与早期的数学摇滚有所不同，早期作品常被认为更具不协和性质。我在下面链接了一篇由 Fecking Bahamas 的 Nickk Hunter 撰写的文章，比我在这里解释得更清楚。}

\eng{Of course, this is not to say that every current math rock band and every song they write is more melodic than their predecessors. Even though it is a niche style of music there are many outliers to any claim that I could make.}

\chn{当然，这并不是说每一支当前的数学摇滚乐队和他们写的每首歌都比前辈更具旋律性。尽管这是一个小众音乐风格，但对于我能提出的任何断言，都存在很多例外。}

\eng{I will begin by explaining major and minor keys, what notes they consist of, and how you can create chords from these notes. Most importantly, I will demonstrate why understanding this basic concept can greatly impact your ability to generate song ideas.}

\chn{我将从解释大调和小调开始，它们由哪些音符组成，以及如何从这些音符构建和弦。最重要的是，我会说明为什么理解这个基本概念会极大地影响你生成歌曲创意的能力。}

\paragraph{Key Signatures}
\eng{A key signature consists of seven notes. Whether it is a major or minor key, a specific set of notes is used, which determines the sharps (\#) or flats (b) for that key. The circle of fifths below shows the sharps and flats for each major and minor key.}

\chn{调号由七个音符组成。无论是大调还是小调，它都会使用一组特定的音符，这决定了该调的升号（\#）或降号（b）。下面的五度圈显示了每个大调和小调的升号与降号。}

\eng{The outside of the circle (red) shows major keys, and the inside (green) shows minor keys. Notice how each major key has a relative minor key vice versa. For example, C major and A minor are relative of each other, meaning they share the same accidentals (which is no sharps or flats), same notes, and chords, but they start from different notes.}

\chn{五度圈外圈（红色）表示大调，内圈（绿色）表示小调。注意每个大调都有一个关系小调，反之亦然。例如 C 大调和 A 小调是关系调，它们共享相同的升降号（此处为无升无降）、相同的音符和和弦，但起始音不同。}

\paragraph{What Notes Are Available and How Many?}
\eng{When counting in sharps, we have these twelve notes: A, A\#, B, C, C\#, D, D\#, E, F, F\#, G, and G\#. And when counting in flats, we have these twelve notes: A, Bb, B, C, Db, D, Eb, E, F, Gb, G, and Ab. It is important to note that A\# and Bb (and others like C\# and Db, D\# and Eb) are the same notes. These are called enharmonic equivalents, and which one you use (sharps or flats) depends on your choice of key signature. A key signature will use seven of the twelve available notes. For example, C major uses C D E F G A B and back to C (the octave above).}

\chn{按照升号计数，我们有这十二个音：A、A\#、B、C、C\#、D、D\#、E、F、F\#、G 和 G\#。按照降号计数，我们有这十二个音：A、Bb、B、C、Db、D、Eb、E、F、Gb、G 和 Ab。需要注意的是，A\# 和 Bb（以及 C\# 和 Db、D\# 和 Eb）实际上是相同的音。这称为等音，使用哪一个（升号还是降号）取决于你选择的调号。一个调号会使用这十二个音中的七个。例如 C 大调使用 C D E F G A B，然后回到高八度的 C。}

\eng{Do not worry about learning all the keys right now, you will pick them up the more you use them, and you can always refer to the circle of fifths to know what sharps or flats are in the key you are using. To add clarity, I recommend watching this video about keys and the circle of fifths.}

\chn{不要担心现在就学会所有调。你用得越多，就会越熟练，而且你随时可以参考五度圈，了解所用调中的升号或降号。为便于理解，我建议观看一段关于调和五度圈的视频。}

\eng{Lastly, you may be asking 'how do I know which key I am in when writing or playing a piece of music?' I would like you to have a little think about how you could work this out, but I will show you a method soon!}

\chn{最后，你可能会问“写歌或演奏时，如何判断自己处于哪个调？”我希望你先思考一下如何解决这个问题，但我很快会展示一种方法！}

\paragraph{Why Is It Useful to Know Key Signatures?}
\eng{There are many benefits to knowing what key we are using when writing or playing songs. To start with, it is an excellent tool for knowing what chords, notes, and scales/arpeggios are available to us.}

\chn{知道写歌或演奏时所用的调有许多好处。首先，它是了解我们可以使用哪些和弦、音符和音阶/琶音的极佳工具。}

\eng{To elaborate, in the key of C major, we have these notes:}

\chn{具体来说，在 C 大调中，我们有以下音：}

\begin{center}
\eng{1 2 3 4 5 6 7 8 (1)}
\chn{1 2 3 4 5 6 7 8 (1)}
\eng{C D E F G A B C}
\chn{C D E F G A B C}
\end{center}

\eng{First off, if you consult the circle of fifths chart above and find C major, you can see there are no sharps or flats (accidentals) in the key of C major. Now you know the appropriate notes to use for a song or an idea in the key of C major (you could play any of the above notes and not hit a bum one). However, careful selection of notes based on music theory can tell us what notes will 'fit' nicely along with chords we choose, and which notes will not (this is not to say there are only correct and incorrect notes to use for a melody, but rather we will get different results from choosing different notes).}

\chn{首先，如果你查看上方的五度圈并找到 C 大调，你会发现 C 大调没有升号或降号（临时记号）。现在你知道了在 C 大调中写歌或构思时可以使用哪些音（你可以弹奏上述任何一个音，不会出现难听的音）。然而，基于乐理的精心音符选择可以告诉我们哪些音更适合与所选和弦搭配，哪些音则不太适合（这并不是说旋律只有正确或错误的音，而是说选择不同的音会产生不同的效果）。}

\paragraph{How Basic Chords Work}
\eng{By knowing the notes in a key signature, we can start to build chords. (For the sake of simplicity, we will stick to standard tuning).}

\chn{通过了解调号中的音符，我们就可以开始构建和弦。（为简化起见，我们仍使用标准调弦。）}

\eng{Continuing with the key of C major, let's build some basic triads (basic chords – those open chord shapes you first learn when learning guitar). We have the chords:}

\chn{继续以 C 大调为例，构建一些基本三和弦（基本和弦——你学吉他时最先学到的那些开放和弦）。我们得到以下和弦：}

\begin{center}
\eng{C D E F G A B}
\chn{C D E F G A B}
\end{center}

\eng{Chords built from a major key follow this order: Major, Minor, Minor, Major, Major, Minor, Diminished. These are often referred to by roman numerals (see above). Upper case for major, and lower case for minor or diminished; I ii iii IV V vi vii°.}

\chn{由大调构建的和弦遵循这个顺序：大和弦、小和弦、小和弦、大和弦、大和弦、小和弦、减和弦。通常用罗马数字表示（见上文）。大写表示大和弦，小写表示小和弦或减和弦；I ii iii IV V vi vii°。}

\eng{Chords built in a minor key follow this order: Minor, Diminished, Major, Minor, Minor, Major, and Major or i ii° III iv v VI VII}

\chn{小调构建的和弦遵循这个顺序：小和弦、减和弦、大和弦、小和弦、小和弦、大和弦、大和弦，表示为 i ii° III iv v VI VII。}

\eng{Here is the key of A minor, which is the relative minor key of the key of C major (same chords but in a different order):}

\chn{下面是 A 小调，它是 C 大调的关系小调（和弦相同但顺序不同）：}

\begin{center}
\eng{A B C D E F G}
\chn{A B C D E F G}
\end{center}

\eng{A chord needs to use at least three notes to be considered a chord, and the distances between the notes in a chord (intervals = spaces) makes it they type of chord it is. (Two notes played at the same time are referred to as intervals, but we could outline a chord with just two notes). Basic chords are made up of three notes, and you are probably already familiar with open chords such as C major:}

\chn{一个和弦至少需要包含三个音，和弦中音符之间的距离（音程＝间隔）决定了它的类型。（同时演奏两个音称为音程，但我们也可以用两个音来概述一个和弦。）基本和弦由三个音组成，你可能已经熟悉像 C 大调这样的开放和弦：}

\eng{The spelling for basic major chords is 1, 3, 5 - (Root, Major 3rd and Perfect 5th are the intervals counting from the C, the root note).}

\chn{基本大三和弦的构成是 1、3、5（从根音 C 开始计数为根音、大三度和纯五度）。}

\eng{If we refer to the C major scale, we can see where the notes that make up a C major chord come from. I have highlighted them in red.}

\chn{如果我们参考 C 大调音阶，就可以看到组成 C 大和弦的音符来自哪里。我已将它们用红色标出。}

\eng{The same applies for basic minor and diminished chords but their spellings are slightly different:}

\chn{对于基本小三和弦和减三和弦也是同样的道理，但它们的构成稍有不同：}

\begin{itemize}[leftmargin=*]
  \item \eng{Major: 1–3–5}
  \chn{大和弦：1–3–5}
  \item \eng{Minor: 1–b3–5 (flat 3rd)}
  \chn{小和弦：1–b3–5（降三度）}
  \item \eng{Diminished: 1–b3–b5 (flat 3rd, flat 5th)}
  \chn{减和弦：1–b3–b5（降三度、降五度）}
\end{itemize}

\eng{You can see this difference in action (the moving third degree and/or fifth degree) below:}

\chn{你可以在下面看到这个差异的实际效果（移动的三度和/或五度）。}

\eng{In a G major chord, there is a major 3rd (B), and when we change to the G minor chord: a flattened 3rd degree, a Bb. Lastly, in the diminished chord we have a flattened 3rd and a flattened 5th Db, and Bb. These intervallic differences are what makes these chords sound the way they do.}

\chn{在 G 大和弦中，有一个大三度（B）；当我们转换到 G 小和弦时，三度变为降三度 Bb。最后，在减和弦中，我们有降三度和降五度 Db、Bb。这些音程差异就是让这些和弦听起来各不相同的原因。}

\paragraph{How Do I Know Which Key Signature Is Used?}
\eng{Here is a method that I learned during my time studying music at university:}

\chn{这是我在大学学习音乐时学到的一种方法：}

\eng{If a key signature is not given, you can often figure out the key by analysing the harmonic elements such as chords, melody, and how different melodic instruments interact.}

\chn{如果没有给出调号，通常可以通过分析和声元素来判断调，例如和弦、旋律，以及不同旋律乐器之间的互动。}

\eng{For example, if you had an intro section of a song using chords: C major, F major, and G major, we can work out the key from this information. I would deduce that this song is in the key of C major, this is because these three chords are all contained in the key of C major, and because the F major and G major chords are right next to each other this tells me they are most likely to be the fourth and fifth chord of a major key (or the VI VII if we were in A minor):}

\chn{例如，如果一首歌的前奏部分使用了 C 大调、F 大调和 G 大调，我们就可以从这些信息判断出这首歌的调。我会推测这首歌是 C 大调，因为这三个和弦都属于 C 大调，而且 F 大调和 G 大调紧挨着，这意味着它们很可能分别是大调中的 IV 和 V（如果是在 A 小调中，则是 VI 和 VII）。}

\begin{center}
\eng{C D E F G A B}
\chn{C D E F G A B}
\eng{Major Minor Minor Major Major Minor Dim}
\chn{大 小 小 大 大 小 减}
\eng{I ii iii IV V vi vii°}
\end{center}

\eng{You can apply the same method if you see two minor chords together or two minor chords that occur together but are spaced out in the chord progression. For example, a C major – D minor – F major – E minor progression (I – ii – IV – iii), the D minor and E minor are spaced out in the progression, but they are next to each other alphabetically in the key of C major.}

\chn{如果你看到两个小和弦并列出现，或者两个小和弦在和弦进行中分布较开，也可以应用同样的方法。例如 C 大调 – D 小调 – F 大调 – E 小调的进行（I – ii – IV – iii），其中 Dm 和 Em 在进行中虽然分开，但在 C 大调的音阶中它们是相邻的。}

\eng{Of course, this method is not bulletproof, especially when you have ambiguous information or less information to work with. For example, let's say you have just two chords: C major and G major, this could be either the key of C major or G major:}

\chn{当然，这个方法并不绝对，尤其是当你得到的信息模糊或不足时。例如，假设你只有两个和弦：C 大调和 G 大调，这可能是 C 大调或 G 大调。}

\begin{center}
\eng{C D E F G A B}
\chn{C D E F G A B}
\eng{Major Minor Minor Major Major Minor Dim}
\chn{大 小 小 大 大 小 减}
\eng{I ii iii IV V vi vii°}
\chn{I ii iii IV V vi vii°}
\eng{G A B C D E F\#}
\chn{G A B C D E F\#}
\eng{Major Minor Minor Major Major Minor Dim}
\chn{大 小 小 大 大 小 减}
\eng{I ii iii IV V vi vii°}
\end{center}

\eng{In cases such as these we generally need more information or take extra steps to work it out. First, I would check if there was an F\# in the melody or a riff, this would tell us it is likely to be G major because that key contains an F\#. Another way you can tell is by extending the chords to sevenths, you would either get a C major 7 and G dominant 7 (extended chords elaborated on in the following section) in the key of C major, or a G major 7 and C major 7 in the key of G major. I would try out both sets of chords and hear which set of chords fits the song better. From this, you can better deduce whether it was the key of G or C major. Of course, it could be either key, and it could be your decision to choose one or even make use of both!}

\chn{在这种情况下，通常需要更多信息或进一步分析。首先，我会检查旋律或即兴片段中是否出现了 F\#，如果有，那很可能是 G 大调，因为该调包含 F\#。另一种判断方法是将和弦扩展到七和弦：在 C 大调中，你会得到 Cmaj7 和 G7（下节将详细讨论扩展和弦）；在 G 大调中，你会得到 Gmaj7 和 Cmaj7。我会尝试这两组和弦，听听哪一组更适合这首歌。由此，你可以更好地判断它是 G 调还是 C 调。当然，它也可能两者皆是，甚至你可以选择混用。}

\eng{It is also worth keeping in mind that even though a song may be situated within a particular key, there may be chords and notes borrowed from other key signatures. There are methods we can use to help identify borrowed chords. And you will see how this is usually done in the borrowed chord segment towards the end of this section.}

\chn{还应记住，尽管一首歌可能位于某个特定调中，但它可能会借用其他调的和弦和音符。我们有一些方法可以帮助识别借用和弦。这一点你将在本节末尾的借用和弦部分看到。}

\paragraph{Task}
\eng{Below, I have given few chord progressions, and I would like you to work out the key of each one. The answers are provided at the end of this section. And keep in mind that some answers may be ambiguous or use more than one key. Fun times! Good luck!}

\chn{下面我给出一些和弦进行，请你判断每个进行所处的调。答案会在本节末提供。请记住，有些答案可能会存在歧义，或者涉及多个调。祝你玩得开心！好运！}

\begin{enumerate}[leftmargin=*]
  \item \eng{E major – A major – G\# minor – F\# minor – E major}
  \chn{E 大调 – A 大调 – G\# 小调 – F\# 小调 – E 大调}
  \item \eng{D minor – G7 – C major}
  \chn{D 小调 – G7 – C 大调}
  \item \eng{D major – E minor – F\# diminished – G major}
  \chn{D 大调 – E 小调 – F\# 减和弦 – G 大调}
  \item \eng{C major – G major – A minor – E minor}
  \chn{C 大调 – G 大调 – A 小调 – E 小调}
  \item \eng{G major – B minor – C major – C minor}
  \chn{G 大调 – B 小调 – C 大调 – C 小调}
\end{enumerate}
